\section{The principal regular insertion after one step}
\label{sec:regular-principal-producer}
We combine the reference attenuation with the comparison to the actual
principal cluster. This supplies the remaining ordinary principal term in
the zero-coefficient transition of Section~\ref{sec:leaf-transfer}.
The forced and nonprincipal terms retain the prices proved in
Section~\ref{sec:projection-payment-producer}.

Let $f$ be a real regular fine insertion on $X$, of price $\tau\ge0$,
and let $Y$ be its selected coarse enlargement. Write $q=C_{B_f}(X)$,
$q'=C_{B_c}(Y)$, and $m_f=e^{a_fq}\tau$ for its fine mass. Its
polynomial moments through degree $d$, in the sense of
\eqref{eq:polynomial-moment-null}, are assumed zero whenever $X$ is
nonwrapping. For the selected transition, assume $X$ has an anchor,
$B_f\le 2B_c$, and the coarse period side $c$ satisfies
$100(1+B_c)\le c$. These are the radius and room conditions used in
the short-support geometry. The order and weight requirements are
\[
 \max\{2,d+1\}\le n_c,\qquad d\le p_f,
 \qquad d+2\le J_f,
\]
where $J_f$ is the fine kernel's available derivative level. Together
with the local regularity required in the preceding sections, these
permit the strip derivatives and the extra insertion derivative used
in the Taylor argument. Let $e=a_c+2$ be the exponent charged by the
coarse reserved mass, and put
\[
 \mathfrak m_c(h)=e^{eq'}h,
 \qquad
 \gamma_L=\frac{L^{-2/5}}{h_d},
 \qquad w_L=\frac{\gamma_L}{2L^d},
\]
where $h_d\ge1$ is the fixed projection and frame cost. We prove a
coarse derivative bound of reserved mass at most $w_Lm_f$ for the
actual principal real insertion.

The constants in the short and long reference estimates, say $A_d$
and $A'_d$, are chosen before $L$. Choose $L$ so that the support
compression inequality holds and
\[
 A_dL^{-d-1}\le w_L/2,
 \qquad A'_dL^{-d-1}\le w_L/2.
\]
This is possible since each reduces to
$4h_d A_d\le L^{3/5}$ or the analogous inequality for $A'_d$.
The strip-to-polynomial bound is chosen before these two constants;
it is uniform in all admissible periods. Only after these choices do we increase
the entrance coupling to meet the reference smallness conditions.

Denote the actual principal insertion by $F$ and the reference
Gaussian integral by $H$. For $q'\le B_c$,
Section~\ref{sec:principal-gaussian-comparison} gives, uniformly over
the forests and replica choices,
\[
 |D^j(F-H)|\le\tau z_\beta,
                       \qquad 0\le j\le n_c,
\]
for the admissible coarse ordered directions. Here $z_\beta\to0$
after $L$ is fixed. Averaging over the forests and replica choices
preserves this bound: the norm of a finite average of differences is
at most their common bound. The enlargement estimate gives
\[
 \mathfrak m_c(\tau z_\beta)
       \le e^{e t_r}z_\beta m_f,
\]
where $t_r$ is the fixed allocation cost. Increase the entrance
coupling so that $e^{e t_r}z_\beta\le w_L/2$.

If $q^3<L$, the stated radius, room and anchor conditions make $X$
nonwrapping by the short-support geometry, so
its assumed moments vanish. The short reference estimate then gives
\[
                    \mathfrak m_c(\text{price of }H)
                         \le A_dL^{-d-1}m_f.
\]
If $q^3\ge L$, use the separate-gauge reference bound. Multiplying
it by $[(1+q)^3/L]^{d+1}\ge1$ inserts $L^{-d-1}$; the extra support
power is absorbed by the unused exponential fine mass, as proved in
Section~\ref{sec:principal-reference-attenuation}. Thus its reserved
mass is at most $A'_dL^{-d-1}m_f$. In either case the reference
price and comparison price add to at most $w_Lm_f$.

It remains to consider $q'>B_c$. Here the principal reserve from
Section~\ref{sec:projection-payment-producer} suffices directly.
The large-support derivative bound has a fixed polynomial cost in
$q'$ multiplied by an exponential with at least one spare unit of
support exponent. For a fixed integer $k$,
\[
 (1+q')^k e^{-q'}
 \le \left[\sup_{t\ge0}(1+t)^k e^{-t/2}\right]e^{-B_c/2}.
\]
The supremum is finite: differentiation of
$k\log(1+t)-t/2$ puts its maximum at $t=\max\{0,2k-1\}$.
The remaining multiplier also contains the covariance and chart
price divided by the Gaussian scale: denote it by $E_L(\beta)$.
It has finite upper power exponent in the incoming coupling, as
proved in Section~\ref{sec:projection-payment-producer}; it is
not a fixed constant. The principal reserved mass is at most
$E_L(\beta)e^{-B_c/2}m_f$. For some fixed $a,C$, eventually
$E_L(\beta)\le C\beta^a$. Since $B_c$ grows quadratically
in $\log\beta_c$, and $\beta_f=\beta_c+$ a fixed shift
after $L$ is chosen,
$\log(C\beta^a e^{-B_c/2})\to-\infty$. Increase the coupling
until $E_L(\beta)e^{-B_c/2}\le w_L$. This estimate applies without a logarithmic
chart on the whole support.

All cases give the asserted principal bound. For a complex insertion,
apply the result to its real and imaginary parts. Each has the same
admissible fine price, and combining their two coarse bounds costs
$2w_Lm_f$. Physical normalization multiplies by $L^d$, giving
\[
                         L^d(2w_Lm_f)=\gamma_Lm_f.
\]
This is the principal share used by the ordinary projection estimate.
It is produced for the same insertion, support, forest average and
transition as the other terms; no replacement of the transition
operator is involved. The case of a primary coefficient of positive
grade uses moments only through one less than that grade and is
handled separately below.
